\section{Definition du probleme}

\subsection{Contexte et motivation}

Aujourd'hui, les contenus d'actualite sont produits en continu et a grande vitesse par un grand nombre de sources numeriques. Les utilisateurs veulent suivre des nouvelles venant de differents pays, de differentes langues et de differents domaines. Cependant, la longueur des articles, le nombre eleve de sources, la repetition du meme evenement sous plusieurs formes et la difficulte d'acceder aux contenus en langues etrangeres rendent l'acces rapide a l'information plus difficile. Cette situation augmente le besoin de systemes capables de presenter, dans les flux d'actualites numeriques, des resumes courts, comprehensibles, fiables et lies a leur source originale.

\subsection{Importance du probleme}

Ce besoin est aussi confirme par les donnees recentes sur la consommation de l'information. Le \textit{Digital News Report 2025}, publie par le Reuters Institute, s'appuie sur des donnees provenant de six continents et de 48 marches. Il montre que la consommation des actualites devient de plus en plus fragmentee, centree sur les plateformes et marquee par un probleme de confiance. Selon ce rapport, \%58 des participants se disent inquiets de ne pas pouvoir distinguer ce qui est vrai et ce qui est faux dans les informations en ligne. Le meme rapport indique aussi que les utilisateurs voient certains usages de l'intelligence artificielle comme utiles pour rendre l'actualite plus accessible: \%27 des participants trouvent utile l'utilisation de l'IA pour resumer les nouvelles, et \%24 pour traduire les nouvelles dans d'autres langues \cite{Newman2025DNR}. De facon similaire, le \textit{Digital News Report 2024} montre qu'environ \%39 des participants evitent parfois ou souvent les nouvelles, et que le sentiment d'etre surcharge par la quantite d'informations a augmente par rapport aux annees precedentes \cite{Newman2024DNR}. Ces resultats montrent que le resume multilingue d'actualites, la fidelite a la source et la presentation rapide de resumes accessibles ne sont pas seulement des problemes techniques, mais aussi des questions liees aux habitudes actuelles de consommation de l'information.

\subsection{Probleme traite par le projet}

Le probleme principal traite par ce projet est le resume automatique de textes d'actualite multilingues, la comparaison objective de differents modeles de resume et l'integration de ces modeles dans une application web fournissant un flux d'actualites en temps reel. Le probleme ne consiste pas seulement a raccourcir un article. Il faut aussi extraire correctement le corps de l'article, rendre le texte utilisable par le modele, produire un resume fidele au texte source, reduire les repetitions et la perte d'information, obtenir des resultats acceptables dans plusieurs langues et permettre au systeme de repondre a l'utilisateur dans un temps raisonnable.

\subsection{Question de recherche}

Dans ce contexte, la question de recherche du projet peut etre formulee ainsi: comment evaluer des modeles de resume d'actualites multilingues selon la qualite, la fidelite a la source, la lisibilite, la couverture de l'information, la vitesse et l'utilisabilite pratique, et comment appliquer ces modeles dans un flux d'actualites provenant de sources reelles? Cette question possede deux dimensions principales: le developpement des modeles et la conception du systeme. Du cote du developpement des modeles, le projet analyse les performances des modeles bases sur mBART entraines dans le cadre du projet et les compare avec un modele mT5 pret a l'emploi, deja entraine dans une etude similaire. Du cote de la conception du systeme, le projet traite la collecte des actualites depuis des sources RSS, l'extraction du texte des articles, la generation des resumes, leur stockage dans une base de donnees et leur presentation dans une interface web.

\subsection{Defis d'evaluation et d'integration}

Un aspect important du probleme est que la qualite d'un resume ne peut pas etre mesuree par une seule metrique. Les metriques classiques comme ROUGE, BLEU et METEOR sont utiles pour mesurer le recouvrement lexical ou phrastique avec un resume de reference. Cependant, dans le contexte du resume d'actualites, elles ne representent pas suffisamment des dimensions comme la fidelite au texte source, le taux de repetition, la clarte, la pertinence, la qualite de la compression et le temps de generation. Pour cette raison, le projet utilise, en plus des metriques classiques, des scores proxy representant des dimensions comme la coherence, l'exactitude, la clarte, la pertinence, l'efficacite, la fidelite a la source et la completude du contenu. Ainsi, le probleme n'est pas traite seulement avec la question ``quel modele obtient le meilleur score classique?'', mais aussi avec la question ``quel modele produit les resumes les plus adaptes a un scenario d'utilisation reel?''.

Un autre aspect important du probleme est la difference entre les jeux de donnees academiques et les flux d'actualites reels. Les jeux de donnees d'evaluation sont generalement composes de textes propres et bien structures, alors que les donnees provenant de sources d'actualites reelles peuvent contenir des structures HTML variees, des contenus incomplets, des textes publicitaires ou de navigation, des repetitions de titres et des problemes de format propres a chaque source. Le probleme du projet ne se limite donc pas a mesurer la performance des modeles; il inclut aussi la mise en place d'une chaine de traitement capable de fonctionner avec des textes bruites provenant de sources reelles. Dans ce cadre, l'extraction du corps des articles, le nettoyage du texte, la generation des resumes, les mecanismes de repli en cas d'echec et la conservation du lien vers la source sont des elements essentiels du systeme.

\subsection{Synthese du probleme}

En conclusion, ce projet traite a la fois le probleme du choix et de l'evaluation des modeles dans le resume d'actualites multilingue, et le probleme de leur integration dans une application en temps reel. La solution attendue est un prototype de bout en bout capable de collecter des actualites dans plusieurs langues, de les resumer, d'evaluer les sorties des modeles avec des metriques multidimensionnelles et de les presenter a l'utilisateur dans une interface web filtrable. Cette approche relie l'evaluation traditionnelle sur jeux de donnees fixes avec les exigences d'un flux d'actualites reel et de l'utilisabilite applicative, ce qui constitue le coeur du probleme traite par le projet.
